% =============================================================================
% Gemeinsame Musterfolien für alle Logovarianten: eine Folie mit Diagrammen,
% eine mit Aufgabenbox und eine mit bewusst langem Titel (prüft, ob Titel und
% Logo kollidieren). Inhalt aus bewegungsgleichungen-wiederholung-folien.tex.
% =============================================================================
\begin{frame}{Steigung im Diagramm}
  \small
  \begin{columns}[T]
    \begin{column}{0.48\textwidth}
      \centering
      \begin{tikzpicture}[scale=0.8]
        \draw[->] (-0.2,0) -- (5,0) node[right] {$t$};
        \draw[->] (0,-0.2) -- (0,3.6) node[above] {$s$};
        \draw[very thick,titleblue] (0.3,0.4) -- (4.5,3.1);
        \draw[thick,groupteal] (1.5,1.171) -- (3.5,1.171);
        \draw[thick,accentorange] (3.5,1.171) -- (3.5,2.457);
        \node[below,groupteal] at (2.5,1.171) {\footnotesize $\Delta t$};
        \node[right,accentorange] at (3.5,1.814) {\footnotesize $\Delta s$};
      \end{tikzpicture}\\[0.3em]
      $s$-$t$-Diagramm:\\
      Steigung $=\dfrac{\Delta s}{\Delta t}=v$
    \end{column}
    \begin{column}{0.48\textwidth}
      \centering
      \begin{tikzpicture}[scale=0.8]
        \draw[->] (-0.2,0) -- (5,0) node[right] {$t$};
        \draw[->] (0,-0.2) -- (0,3.6) node[above] {$v$};
        \draw[very thick,accentorange] (0.3,0.6) -- (4.5,3.1);
        \draw[thick,groupteal] (1.5,1.314) -- (3.5,1.314);
        \draw[thick,accentpurple] (3.5,1.314) -- (3.5,2.505);
        \node[below,groupteal] at (2.5,1.314) {\footnotesize $\Delta t$};
        \node[right,accentpurple] at (3.5,1.91) {\footnotesize $\Delta v$};
      \end{tikzpicture}\\[0.3em]
      $v$-$t$-Diagramm:\\
      Steigung $=\dfrac{\Delta v}{\Delta t}=a$
    \end{column}
  \end{columns}
\end{frame}

\begin{frame}{Aufgabe 4: Gleichung 2 anwenden}
  \begin{aufgabebox}[Aufgabe 4]
    \small
    Ein Auto fährt auf einem Autobahnbeschleunigungsstreifen. Eine
    Stoppuhr läuft schon. Bei $t_0=\SI{3}{\second}$ hat das Auto
    $v_0=\SI{15}{\meter\per\second}$. Ab da beschleunigt es konstant mit
    $a=\SI{2}{\meter\per\second\squared}$.
    \begin{enumerate}
      \item[a)] Welche Geschwindigkeit hat das Auto, wenn die Stoppuhr
        $t=\SI{8}{\second}$ zeigt? Gib sie auch in
        $\si{\kilo\meter\per\hour}$ an.
      \item[b)] Welche Zeit zeigt die Stoppuhr, wenn das Auto
        $\SI{27}{\meter\per\second}$ erreicht?
    \end{enumerate}
  \end{aufgabebox}
\end{frame}

\begin{frame}{Gleichung 2: Geschwindigkeit ($a$ konstant)}
  \small
  \begin{itemize}
    \item Bekannt: $a=\dfrac{\Delta v}{\Delta t}$
    \item Mal $\Delta t$ und $\Delta$ ausschreiben:
      $a\cdot(t-t_0)=v-v_0$
    \item Plus $v_0$: $\quad v(t)=v_0+a\cdot(t-t_0)$
  \end{itemize}
  \vspace{0.5em}
  Diese Folie hat einen langen Titel, damit man sieht, ob Titel und Logo
  sich in die Quere kommen.
\end{frame}
